\section{Uncertainties and PoC Plan}
\label{sec:poc}

\subsection{Ranked Uncertainty Register}

Table~\ref{tab:uncertainties} ranks the three uncertainties that
would most damage the project if wrong; R1 is attacked first.

\begin{table*}[t]
\centering
\footnotesize
\setlength{\tabcolsep}{3pt}
\caption{Three ranked uncertainties. R1 is attacked first, since it would invalidate the most work if wrong.}
\label{tab:uncertainties}
\begin{tabular}{p{1.9cm}p{4.95cm}p{4.95cm}p{4.95cm}}
\toprule
\textbf{Field} & \textbf{R1 (ranked 1st)} & \textbf{R2 (ranked 2nd)} & \textbf{R3 (ranked 3rd)} \\
\midrule
Statement &
The fused classifier may not reach the 15\% error target on
wearers it was not trained on. &
Delivery under a genuinely degraded link may exceed the 30\,s
target, or messages may be lost rather than retried. &
Fusing chest and/or thigh features with the wrist may not lower
error compared to wrist alone, undermining the reason for the
custom board. \\[0.6em]

Assumption &
Motion features learned from five volunteers generalize to a
sixth, unseen wearer. &
The report/retry logic behaves correctly under injected delay
and packet loss, not only under a healthy link. &
Spatial diversity across sensing sites decorrelates misleading
readings enough to matter for these three activities. \\[0.6em]

Consequence &
REQ-AC-02 unmet; the AI claim in the PoC question fails. &
REQ-AC-04/05 unmet; the IoT claim in the PoC question fails. &
The custom board's contribution goes unproven, though the
system still works via the wrist-only fallback tier. \\[0.6em]

Probability &
Unknown with only five volunteers -- treated as an uncertainty,
not a risk. &
Unknown until tested under real degraded conditions --
uncertainty. &
Unknown -- uncertainty. \\[0.6em]

Evidence\newline required &
Leave-one-wearer-out error measured across all five volunteers. &
Delivery latency and message loss measured under the
degraded-link profile and a full link outage (defined below). &
Wrist-only vs.\ board-only vs.\ full-fusion error compared on
the same leave-one-wearer-out evaluation (REQ-AC-09). \\[0.6em]

Mitigation &
If exceeded, revisit the feature set or record additional
volunteers before final submission. &
If exceeded, relax the degraded-link time budget with a
recorded reason, or widen the retry buffer. &
If fusion does not help, report that honestly and treat
wrist-only as the primary design in Stage 1. \\
\bottomrule
\end{tabular}
\end{table*}

\subsection{Proof of Concept Plan}

\textbf{Question:} Does fusing wrist, chest and/or thigh IMU data
improve activity classification over wrist alone, and does the
link deliver each result on time when degraded, with no loss once
a full outage recovers?

\textbf{Link conditions} (applied by the companion computer's
test rig, which carries the traffic but does no classification):
\begin{itemize}
\item \emph{Healthy:} no added delay or loss, watch and board in
the same room as the access point.
\item \emph{Degraded:} 500\,ms added delay and 20\% packet loss
on the watch-to-computer path.
\item \emph{Outage:} link fully cut for 60\,s, then restored.
\end{itemize}

\textbf{Measurements:}
\begin{itemize}
\item Error rate: wrist-only vs.\ board-only (chest and/or thigh)
vs.\ full fusion, per left-out wearer.
\item Activity-change delivery time to the computer: healthy vs.\
degraded link.
\item Messages lost and recovery behaviour during a full link
outage.
\end{itemize}

\textbf{Pass criterion (fixed in advance):} fused mean error at
most 15\% across left-out wearers, reported alongside wrist-only
error; each activity change reaches the computer within 10\,s on
a healthy link and 30\,s on a degraded link; no message lost once
a full outage recovers. Fusion is kept as the primary design only
if its error is lower than wrist-only; otherwise wrist-only
becomes primary (R3 mitigation).
